\begin{figure}[htbp]
\centering
\begin{tikzpicture}
\begin{axis}[width=0.86\linewidth,height=5.8cm,ybar,bar width=15pt,
  ymin=0,ymax=11,ytick={0,2,4,6,8,10},ylabel={Successes (of 10)},
  symbolic x coords={String,Sort,Indexed},xtick=data,
  xticklabels={String concat.,Sort selection,Indexed lookup},
  tick label style={font=\small},nodes near coords,
  every node near coord/.append style={font=\small},
  axis x line*=bottom,axis y line*=left,grid=major,
  grid style={draw=gray!18},enlarge x limits=0.17,
  legend style={at={(0.5,-0.2)},anchor=north,draw=none,
  legend columns=2,font=\small}]
\addplot[fill=passgreen!75,draw=passgreen] coordinates
  {(String,7) (Sort,7) (Indexed,0)};
\addplot[fill=deepblue!70,draw=deepblue] coordinates
  {(String,8) (Sort,8) (Indexed,0)};
\legend{v5,v6.1}
\end{axis}
\end{tikzpicture}
\caption{Family-heldout decomposition. v6.1 gains one string-concatenation
and one sort-selection case; both adapters remain 0/10 on indexed lookup.
Each bar is a count over ten cases, not a population estimate.}
\label{fig:heldoutfamilies}
\end{figure}
